\section{Normal estimation}
\label{sec-normals}

The areas of projections are related to local geometric
information. More precisely, we show in this section that, in the case
where $X$ is a hypersurface (a $d-1$-submanifold of $\R^d$), the
normalized projected areas of dimension $d-1$ are close to the
absolute values of the normal vector components. This induces a
multigrid convergent normal estimator for digitized surfaces (see
Corollary~\ref{cor:cubical-normal-estimation} and the following
remark). From now on, we write $\widehat{i}$ for the $(d-1)$-element subset of
indices $\{1,\ldots,d\} \setminus \{ i \}$.

\begin{theorem}
\label{thm:normal-recovery}
Let $T\subset\mathbb R^d$ be a $(d-1)$-dimensional affine subspace
with unit normal vector $n=(n_1,\ldots,n_d)$, and let $p \in T$.
Then \( V_{\widehat{i}}(T,p,r)=\omega_{d-1}r^{d-1}|n_i|. \)
Consequently, the $d$ $\pi_{\widehat{i}}$-projection volumes uniquely
determine the absolute values of the components
of the unit normal vector:
\[
W_{d-1}(T,p,r) := \sum_{i=1}^d V_{\widehat{i}}(T,p,r)^2
=
\omega_{d-1}^2r^{2(d-1)},
~\text{and}~
|n_i|=\frac{V_{\widehat{i}}(T,p,r)}
{\sqrt{W_{d-1}(T,p,r)}}.
\]
\end{theorem}

\begin{proof}
Let $A\in\mathbb R^{d\times(d-1)}$ have orthonormal columns spanning
the direction space of $T$, so that $A^\top A=\mathrm{Id}_{d-1}$.  For each $i, 1\Le i \Le
d$, let $A_{\widehat i}$ be the square matrix obtained by deleting the
$i$-th row of $A$.  The restriction of $\pi_{\widehat i}$ to $T$ is
represented by $A_{\widehat i}$.  Hence, by standard area computation, 
\(
V_{\widehat{i}}(T,p,r)=\omega_{d-1}r^{d-1}|\det A_{\widehat i}|.
\)

\begin{sloppypar}
Consider the vector of signed maximal minors $c=(c_1,\ldots,c_d)$,
with $c_i = (-1)^{i+1} \det A_{\widehat i}$.  By the cofactor expansion,
$c$ is orthogonal to every column of $A$, and therefore $c$ is
parallel to $n$.
Moreover, applying the Cauchy--Binet formula gives
\(
|c|^2
=
\sum_{i=1}^d\det(A_{\widehat i})^2
=
\det(A^\top A)=1.
\)
Since $n$ is a unit normal vector and $c$ is parallel to $n$, we obtain
$c=\pm n$. In particular,
\(
|\det A_{\widehat i}|=|n_i|.
\)
Substitution into the expression for $V_{\widehat{i}}(T,p,r)$
proves the first identity. The remaining
identities follow from $|n|=1$. \qed
\end{sloppypar}
\end{proof}



%% \begin{corollary}%[Local estimation of the absolute normal components]
%% \label{cor:local-normal-estimation}
%% Let $X\subset\mathbb R^d$ be a $C^2$ hypersurface
%% without boundary, with reach at least $\rho>0$.
%% Let $p\in X$, and let
%% $n(p)=(n_1(p),\ldots,n_d(p))$
%% be a unit normal vector at $p$.
%% Define the $i$-th component of the  normal estimator by
%% \( \widehat n_i(p,r)
%% :=
%% V_{\widehat{i}}(X,p,r) / \sqrt{W_{d-1}(X,p,r)}.
%% \)

%% There exist constants $\alpha_0>0,C>0$, depending
%% only on $d$, such that, whenever
%% $0<r<\alpha_0\rho$, the estimator is well-defined
%% and satisfies
%% \[
%% \left\| \widehat n(p,r)-|n(p)| \right\|_2
%% \le C\frac r\rho,
%% \]
%% where \(|n(p)|=(|n_1(p)|,\ldots,|n_d(p)|). \)
%% \end{corollary}

%% \begin{proof}
%% Let $T=T_pX$ and let $B_T(p,r)$ denote the
%% $(d-1)$-dimensional ball of radius $r$
%% contained in the affine tangent hyperplane
%% $p+T$.

%% By \Equ{eq:stabilite-VJ} of the local stability theorem
%% (\RefTheorem{thm:stabilite-projections}), using $Y=X \cap B(p,r)$ and
%% $\varepsilon=0$ for hypothesis \Equ{eq:encadrement-Y}, there exists a
%% constant $C_d>0$ such that, for every $i$,
%% \[
%% \left|
%% V_{\widehat{i}}(X,p,r)
%% - \mathcal H^{d-1}
%% \bigl(\pi_{\widehat i}(B_T(p,r))\bigr)
%% \right|
%% \le
%% C_d\frac{r^d}{\rho}.
%% \]
%% %% The projection $\pi_{\widehat{i}}$ of a shape belonging to $T$ multiplies its measure by $|n \cdot e_i|=|n_i|$, (specific case of \Equ{eq-vol-proj-ball-1} too), so
%% By \RefTheorem{thm:normal-recovery}, 
%% \(
%% \mathcal H^{d-1}
%% \bigl(\pi_{\widehat i}(B_T(p,r))\bigr)
%% =
%% \omega_{d-1}r^{d-1}|n_i(p)|.
%% \)
%% Consequently, \[\left|V_{\widehat{i}}(X,p,r)-\omega_{d-1}r^{d-1}|n_i(p)|\right| \Le C_d\frac{r^d}{\rho}.\]

%% We set $a=\omega_{d-1}r^{d-1}$,
%% $v=(V_{\widehat{1}}(X,p,r),\ldots,V_{\widehat{d}}(X,p,r))$ and
%% $u=|n(p)|$.  Since $\|u\|_2=1$, the preceding relations translate to
%% $\|v-au\|_2 \Le \sqrt d\,C_d\frac{r^d}{\rho}$.

%% Since \(X\) is a smooth hypersurface, at least one projection is
%% locally nonsingular at \(p\). Hence \(v\ne0\) for every \(r>0\).
%% It follows that:

%% \begin{align*}
%% \left\|\widehat n(p,r)-|n(p)|\right\|_2
%% &= \left\|
%% \frac{v}{\|v\|_2}-u
%% \right\|_2
%% \Le \frac{1}{a}\left(\left\| \frac{a}{\|v\|_2}v - v \right\|_2 + \|v-au\|_2 \right) \\
%% &= \frac{1}{a}\left(|a-\|v\|_2| + \|v-au\|_2 \right) \Le \frac{2}{a}\|v-au\|_2 \Le \frac{2\sqrt d\,C_d}
%% {\omega_{d-1}}
%% \frac r\rho. 
%% \end{align*}
%% This proves the claim. \qed
%% \end{proof}


\begin{corollary}%[Normal estimation from a cubical discretization]
\label{cor:cubical-normal-estimation}
Let $X\subset\mathbb R^d$ be a $C^2$ hypersurface without boundary,
with reach at least $\rho>0$.  Let $p\in X$, and let
$n(p)=(n_1(p),\ldots,n_d(p))$ be a unit normal vector at $p$.  Let
$U_h(X)$ be the $h$-voxelization of $X$.  For each
$i\in\{1,\ldots,d\}$, set the $i$-th component of the normal estimator
$\widehat{n}_{h}(p,r)$ as $\widehat{n}_{i,h}(p,r):=
V_{\widehat{i}}(U_h(X),p,r)/\sqrt{W_{d-1}(U_h(X),p,r)}$.
 
There exist constants $\alpha_0>0,C>0$, depending only on $d$, such
that, whenever $0 < h \Le r$ and $0 < r < \min \{ \alpha_0,
\frac{1}{1+\sqrt{d}} \} \rho$, the estimator is well-defined and
satisfies
\[
\left|\widehat{n}_h(p,r)-(|n_1(p)|,\ldots,|n_d(p)|)|\right|
\Le C \left(\frac{h}{r}+\frac{r}{\rho} \right).
\]
\end{corollary}

\begin{proof}
  Let $T=T_pX$ and let $B_T(p,r)$ denote the $(d-1)$-dimensional ball
  of radius $r$ contained in the affine tangent hyperplane $p+T$. Now
  let $Y=U_h(X) \cap B(p,r)$. For
  $\varepsilon=\left(1+\frac{2}{\sqrt3}\right)\sqrt d\,h$ and the
  required hypothesis on $h$ and $r$,
  \RefTheorem{thm:cubical-admissible-set} holds and $Y$ satisfies the
  bounds of hypothesis \Equ{eq:encadrement-Y} of the stability theorem
  (\RefTheorem{thm:stabilite-projections}).  By its equation
  \Equ{eq:stabilite-VJ}, there exists a constant $C_d>0$ such that,
  for every $i$,
\[
\left|
V_{\widehat{i}}(Y,p,r)
- \mathcal H^{d-1}
\bigl(\pi_{\widehat i}(B_T(p,r))\bigr)
\right|
\Le
C_d r^{d-1}\left(\frac{h}{r}+\frac{r}{\rho}\right).
\]


By \RefTheorem{thm:normal-recovery}, 
\(
\mathcal H^{d-1}
\bigl(\pi_{\widehat i}(B_T(p,r))\bigr)
=
\omega_{d-1}r^{d-1}|n_i(p)|.
\)
Consequently,
\[
\left|V_{\widehat{i}}(Y,p,r)-\omega_{d-1}r^{d-1}|n_i(p)|\right|
\Le C_d r^{d-1}\left(\frac{h}{r}+\frac{r}{\rho}\right).
\]

We set $a=\omega_{d-1}r^{d-1}$,
$v=(V_{\widehat{1}}(Y,p,r),\ldots,V_{\widehat{d}}(Y,p,r))$ and
$u=(|n_1(p)|,$ $\ldots$, $|n_d(p)|)$.  Since $|u|=1$, the preceding relations translate to
$|v-au| \Le \sqrt d\,C_d r^{d-1}\left(\frac{h}{r}+\frac{r}{\rho}\right)$.

Since $X\cap B(p,r)\subseteq Y$ and $X$ is a smooth hypersurface, at
least one coordinate projection is locally nonsingular at $p$. Hence
$v\ne 0$ for every $r>0$.
%% For $r>0$, the vector $v$ is not null since $V_{\widehat{i}}(Y,p,r) > 0$ ($Y$ contains at least the projection of the cube containing $p$).
It follows that:

\begin{align*}
\left|\widehat n_h(p,r)-|n(p)|\right|
&= \left|
\frac{v}{|v|}-u
\right|
\Le \frac{1}{a}\left(\left| \frac{a}{|v|}v - v \right| + |v-au| \right) \\
&= \frac{1}{a}\left(|a-|v|| + |v-au| \right) \Le \frac{2}{a}|v-au|
\Le \frac{2\sqrt d\,C_d}
    {\omega_{d-1}}
    \left(\frac{h}{r}+\frac{r}{\rho}\right).
\end{align*}
This proves the claim. 
%% From this point, the proof is entirely similar to the proof of Corollary~\ref{cor:local-normal-estimation}. The vector $v$ is not null since $V_{\widehat{i}}(Y,p,r) \Ge V_{\widehat{i}}(X,p,r)$, so we conclude that:
%% \[
%% \left\|\widehat n_h(p,r)-|n(p)|\right\|_2
%% \Le \frac{2\sqrt d\,C_d}{\omega_{d-1}}
%% \left(\frac{h}{r}+\frac{r}{\rho}\right).
%% \]
\qed
\end{proof}


\begin{remark}[Sign ambiguity and oriented normals]
The projection volumes determine only the absolute values
of the normal components.
%
For digitized boundaries of solid objects, however, the orientation of
the surfels provides additional information.  In particular, algebraic
sums of surfel areas, weighted by their oriented coordinate normals
$\pm e_i$, may be used to construct an estimator of the signed normal
components.
%% Although these signed quantities should be distinguished
%% from the unsigned projection volumes considered above, they are
%% analogous per space orthant.
Although these quantities are closely related to the local
projection measures, establishing the convergence of the
resulting signed normal estimator is left for future work.
%% Although they share similitudes with the local measures, proving that
%% they lead to a convergent normal estimator remains a perspective to
%% this work.
%
Finally, the error bound suggests balancing the discretization
and curvature contributions by choosing
$r=\Theta(\sqrt{\rho h})$.
\end{remark}

%% \begin{remark}[Sign ambiguity]
%% The projection volumes determine only the absolute values of the
%% normal components.  Although it is not developed here, if the
%% discretization $S_h$ is the digitized boundary of a volume $V$, then
%% the orientations of the surfels in the vicinity of $p$ gives you the
%% sign of each components. Even better, counting surrounding surfels
%% with orthogonal vector $\pm e_i$ as $\pm 1$ gives you the (oriented)
%% measures $V_{\widehat{i}}(S_h,p,r)$, hence the oriented normal
%% $\widehat{n}_h(p,r)$. The error term again suggests an optimal radius
%% of $\sqrt{\rho h}$.
%% %% In particular, the relative signs of
%% %% $n_1(p),\ldots,n_d(p)$ cannot be recovered
%% %% from these measurements alone.
%% %% Additional geometric information is needed
%% %% to reconstruct the full normal direction.
%% \end{remark}
